\clearpage
\fancyhead[L]{\small GAUSS TO VON ZACH.}
\hypertarget{native-286-lower}{}
\begin{center}
\textbf{GAUSS TO VON ZACH.}

\hrulefill

Monatliche Correspondenz. Vol. XV. Pp. 377. 378. April 1807.

\hrulefill
\end{center}

\begin{flushright}
Brunswick, 10 March 1807.
\end{flushright}

\par\indent I have the honor of enclosing my first observations of \emph{Pallas} this year. Bad weather has so far kept me from making further observations.

\begin{center}\begingroup
\fontsize{9}{11}\selectfont
\setlength{\tabcolsep}{3pt}\renewcommand{\arraystretch}{1.15}
\begin{tabular}{l|l|l|l}
1807 & Mean time & \shortstack{Apparent right\\ascension $\PallasSymbol$} & \shortstack{Apparent north\\declination $\PallasSymbol$} \\
\hline
February 28 & $11^{\mathrm u}\,55^{\mathrm m}\,42^{\mathrm s}$ & $236^\circ\;48'\;6''$ & $8^\circ\;10'\;15''$ \\
March \phantom{0}3 & $12\;\;55\;\;25$ & $237\;\;15\;\;5$ & $8\;\;54\;\;19$ \\
--- \phantom{0}5 & $12\;\;22\;\;2$ & $237\;\;30\;\;41$ & $9\;\;24\;\;52$ \\
\end{tabular}\endgroup\end{center}

\par\indent The first two observations are less reliable, especially as regards the declinations. The third, made under more favorable conditions, is as accurate as my equip%
\SourceLeaf{287}%
ment allows. Two ninth-magnitude stars from the \emph{Hist. cél.} served for comparison on the fifth of March. Their apparent positions were taken as follows

\begin{center}\begingroup
\fontsize{9}{11}\selectfont
\setlength{\tabcolsep}{3pt}\renewcommand{\arraystretch}{1.15}
\begin{tabular}{l|l}
Right ascension & North declination \\
\hline
$237^\circ\;19'\;12''$ & $9^\circ\;27'\;50''$ \\
$237\;\;31\;\;13$ & $9\;\;28\;\;8$ \\
\end{tabular}\endgroup\end{center}

\par\indent According to these observations, my ephemeris gives the declinations accurately, but the right ascensions are 12 to 13 minutes too large. That difference still seems to be growing. I am most grateful for the Milan observations of \emph{Ceres} from 1806. They have now given me the precision I wanted for this opposition as well. I do find, however, that the declination on 4 January is 1 minute too small and should be $29^\circ\;21'\;29''3$ (February 1807 issue, \href{https://archive.org/details/bub_gb_sKAAAAAAMAAJ/page/n164/mode/1up}{p.~159}). This is no copying error, since the latitude is also too small, and Mr. \textsc{Carlini} must by chance also have been off by 1 minute in the latitude calculated from the elements. Please bring this to the attention of that skilled astronomer when you have a chance. My latest result agrees almost exactly with the one Mr. \textsc{Carlini} calculated. I find the opposition on 3 January 1806, at $12^{\mathrm u}\,20^{\mathrm m}\,12^{\mathrm s}$ mean time in Milan, at $103^\circ\;1'\;9''4$ true longitude and $6^\circ\;28'\;32''4$ geocentric northern latitude. I will now wait to use this result until this year’s opposition has also been observed.

